\section{Fine-Graining Isometries for Anyonic Lattice Models}
\label{sec:finegraining}

We construct fine-graining isometries $V: \catH_L \to \catH_{2L}$ that embed the $L$-site anyonic Hilbert space into the $2L$-site space, doubling the lattice resolution while preserving the physics. These isometries are the central ingredient of the operator algebraic renormalization (OAR) programme~\cite{OsborneStottmeister2022,Stottmeister2021}, providing the inductive limit construction $\mathcal{A}_\infty = \varinjlim \mathcal{A}_N$ that yields the continuum C*-algebra.

\subsection{Setup}

Let $\catC$ be a fusion category with simple objects $X_1 = \one, X_2, \ldots, X_d$. We consider $N$ hard-core anyons on an $L$-site lattice with open boundary conditions. The Hilbert space decomposes as
\begin{equation}
    \catH_N(L) = \bigoplus_{\substack{\text{positions } x_1 < \cdots < x_N \\ \text{labels } a_1, \ldots, a_N}} \Mor(X_c, X_{a_1} \otimes \cdots \otimes X_{a_N})
\end{equation}
where the direct sum runs over all hard-core configurations and all label assignments from the non-vacuum simples.

\subsection{The one-particle scaling map}

The fine-graining begins with a one-particle isometry $R: \mathbb{C}^L \to \mathbb{C}^{2L}$ built from a compactly supported wavelet scaling function. For the Daubechies-$K$ wavelet with filter coefficients $h_0, \ldots, h_{2K-1}$, the matrix elements are
\begin{equation}
    R_{j,i} = h_{j - 2(i-1)}, \qquad i = 1, \ldots, L, \quad j = 1, \ldots, 2L
\end{equation}
with boundary corrections via L\"owdin (symmetric) orthogonalisation to ensure $R^\dagger R = I_L$ exactly. For the Haar wavelet ($K=1$, $h = [1/\sqrt{2}, 1/\sqrt{2}]$), each coarse site $i$ maps to an equal superposition of fine sites $2i-1$ and $2i$:
\begin{equation}
    R\ket{i} = \frac{1}{\sqrt{2}}\bigl(\ket{2i-1} + \ket{2i}\bigr).
\end{equation}

\subsection{The product map}
\label{sec:product-map}

The simplest many-body lift spreads each particle independently:
\begin{equation}
\label{eq:product-map}
    V_{\mathrm{prod}}\ket{x_1, \ldots, x_N; \vec{a}; \vec{e}; c}
    = \sum_{\substack{y_1 < \cdots < y_N}} \prod_{p=1}^N R_{y_p, x_p} \; \ket{y_1, \ldots, y_N; \vec{a}; \vec{e}; c}
\end{equation}
where labels $\vec{a}$ and fusion-tree intermediates $\vec{e}$ are preserved. This map respects the fusion-tree structure because the hopping amplitude in any fusion category is $+1$ (the F-symbols involving vacuum are trivial).

\begin{proposition}[Product map isometry for Haar]
\label{prop:haar-isometry}
For the Haar wavelet, the product map~\eqref{eq:product-map} satisfies $V_{\mathrm{prod}}^\dagger V_{\mathrm{prod}} = I$ for any fusion category $\catC$.
\end{proposition}

\begin{proof}
For Haar, particle $p$ at coarse site $x_p$ spreads to fine sites $\{2x_p - 1, 2x_p\}$. Since $x_1 < \cdots < x_N$ are distinct, the fine-site supports $\{2x_p - 1, 2x_p\}$ are pairwise disjoint. The constrained sum over $y_1 < \cdots < y_N$ factorises:
\begin{equation}
    (V^\dagger V)_{\alpha\beta} = \delta_{\alpha\beta} \prod_{p=1}^N \sum_{y_p \in \{2x_p - 1, 2x_p\}} |R_{y_p, x_p}|^2 = \delta_{\alpha\beta} \prod_{p=1}^N 1 = \delta_{\alpha\beta}.
\end{equation}
The disjointness of supports is the key: it eliminates the hard-core constraint between particles, allowing the sum to factorise.
\end{proof}

\subsection{Failure of the product map for overlapping wavelets}

For wavelets with wider support (e.g., Daubechies D4 with 4 filter coefficients), adjacent particles at coarse sites $i$ and $i+1$ spread to overlapping fine-site ranges: particle $i$ to $\{2i{-}1, 2i, 2i{+}1, 2i{+}2\}$ and particle $i{+}1$ to $\{2i{+}1, 2i{+}2, 2i{+}3, 2i{+}4\}$. The overlap at $\{2i{+}1, 2i{+}2\}$ prevents the sum from factorising, and $V_{\mathrm{prod}}^\dagger V_{\mathrm{prod}} \neq I$.

\begin{remark}
Numerically, the deficit grows with the number of adjacent particle pairs: for D4 on $L=3$ Fibonacci, $\|V^\dagger V - I\| \approx 0.11$ for the product map.
\end{remark}

\subsection{The Slater determinant lift for identical particles}

For a single non-vacuum simple (e.g., $\tau$ in Fibonacci, $\psi$ in sVec), all particles carry the same label. The Hilbert space factorises as
\begin{equation}
\label{eq:factorisation}
    \catH_N(L) = \Lambda^N(\mathbb{C}^L) \otimes \Mor(X_c, X_a^{\otimes N})
\end{equation}
where $\Lambda^N(\mathbb{C}^L) \cong \mathbb{C}^{\binom{L}{N}}$ is the $N$-th exterior power (the hard-core position space) and $\Mor(X_c, X_a^{\otimes N})$ is the fusion-tree internal space.

\begin{theorem}[Slater lift for identical anyons]
\label{thm:slater}
For a fusion category $\catC$ with a single non-vacuum simple object $X_a$, the fine-graining isometry
\begin{equation}
    V = \Lambda^N(R) \otimes I_{\Mor(X_c, X_a^{\otimes N})}
\end{equation}
satisfies $V^\dagger V = I$ for \emph{any} one-particle isometry $R: \mathbb{C}^L \to \mathbb{C}^{2L}$.
\end{theorem}

\begin{proof}
The position part $\Lambda^N(R)$ has matrix elements $(\Lambda^N(R))_{S,T} = \det(R_{S,T})$ for $N$-element subsets $S \subset \{1,\ldots,2L\}$ and $T \subset \{1,\ldots,L\}$. By Cauchy--Binet:
\begin{equation}
    (\Lambda^N(R)^\dagger \Lambda^N(R))_{T,T'} = \sum_{|S|=N} \det(R_{S,T})^* \det(R_{S,T'}) = \det(R^\dagger R)_{T,T'} = \det(I_L)_{T,T'} = \delta_{TT'}.
\end{equation}
Since the internal part is tensored with the identity, $V^\dagger V = I \otimes I = I$.
\end{proof}

\begin{remark}
The antisymmetric signs in $\det$ are a \emph{mathematical} necessity (for Cauchy--Binet), not a statement about fermionic statistics. The same construction works for Fibonacci ($\tau$ anyons with non-abelian braiding), sVec ($\psi$ fermions), or any category with a single particle type.
\end{remark}

\subsection{The multi-species problem}

For categories with multiple non-vacuum simples (e.g., Ising with $\psi$ and $\sigma$), configurations can have mixed labels. The factorisation~\eqref{eq:factorisation} \emph{breaks}: the fusion-tree space $\Mor(X_c, X_{a_1} \otimes \cdots \otimes X_{a_N})$ depends on the specific label sequence $\vec{a}$, which is tied to the positions.

The Slater determinant lift fails for mixed labels because:
\begin{enumerate}
    \item The determinant introduces cross-terms where particles are permuted.
    \item Permuting particles of different types changes the label sequence $\vec{a}$.
    \item Different label sequences correspond to different fusion trees (different bra states).
    \item The cross-terms contribute to \emph{different} fine states, preventing Cauchy--Binet cancellation.
\end{enumerate}

We investigated several alternative constructions:
\begin{itemize}
    \item \textbf{Categorical determinant} $\sum_{\sigma \in S_N} B(\sigma) \prod_p R_{y_p, x_{\sigma(p)}}$ with braiding weights: reduces the error by $\sim$50\% but is not isometric (Cauchy--Binet requires scalar phases, not matrix-valued braiding).
    \item \textbf{Species-enlarged exterior algebra} $\Lambda^N(\mathbb{C}^{Lk})$: the hard-core-on-sites projection (not just hard-core-on-$({\rm site},{\rm species})$ pairs) breaks Cauchy--Binet.
    \item \textbf{Label-grouped Slater}: Slater within same-label groups, product across groups. Partial fix only (handles same-species overlaps but not cross-species).
\end{itemize}

\subsection{The normalized product map: the universal construction}
\label{sec:normalized-product}

\begin{theorem}[Universal fine-graining isometry]
\label{thm:normalized-product}
For any fusion category $\catC$ with any number of particle types, and any one-particle isometry $R$, the \emph{normalized product map}
\begin{equation}
\label{eq:normalized-product}
    V_{\mathrm{iso}} = V_{\mathrm{prod}} \cdot \operatorname{diag}\!\left(\frac{1}{g(\vec{x})}\right)
\end{equation}
is isometric ($V_{\mathrm{iso}}^\dagger V_{\mathrm{iso}} = I$), where $g(\vec{x}) = \|V_{\mathrm{prod}}\ket{\vec{x}, \vec{a}, \vec{e}, c}\|$ is the column norm of the product map.
\end{theorem}

This relies on two structural properties of the product map, which we have verified numerically for sVec, Fibonacci, and Ising with both Haar and D4 wavelets for $L = 3, 4, 5$:

\begin{claim}[Diagonal Gram matrix]
\label{claim:diagonal-G}
The Gram matrix $G = V_{\mathrm{prod}}^\dagger V_{\mathrm{prod}}$ is diagonal: different coarse states map to orthogonal fine-state superpositions.
\end{claim}

\begin{claim}[Position-only normalization]
\label{claim:position-norm}
The column norm $g(\vec{x}, \vec{a}, \vec{e}, c) = g(\vec{x})$ depends only on the particle positions $\vec{x}$, not on the labels $\vec{a}$, intermediates $\vec{e}$, or total charge $c$. In particular, it is \emph{category-independent}.
\end{claim}

The normalization factor has a transparent geometric interpretation:
\begin{equation}
    g(\vec{x})^2 = \sum_{y_1 < \cdots < y_N} \prod_{p=1}^N R_{y_p, x_p}^2
\end{equation}
which is the probability that $N$ independently spread particles all land at distinct sites.

\begin{remark}
For Haar wavelets, $g(\vec{x}) = 1$ for all configurations (non-overlapping supports), recovering the product map. For D4, $g < 1$ when adjacent coarse sites have overlapping fine-site supports. The correction is purely geometric and does not involve any categorical data (F-symbols, R-symbols, etc.).
\end{remark}

\subsection{Number-changing fine-graining isometry}
\label{sec:number-changing}

The normalized product map~\eqref{eq:normalized-product} fixes the deficit $D = I - V_{\mathrm{prod}}^\dagger V_{\mathrm{prod}}$ by rescaling columns, which is mathematically clean but physically artificial. We now show that the deficit can be filled \emph{physically} by pair-creation terms, yielding a number-changing isometry
\begin{equation}
\label{eq:number-changing}
    V = V_0 + V_2, \qquad V_0: \catH_N(L) \to \catH_N(2L), \quad V_2: \catH_N(L) \to \catH_{N+2}(2L)
\end{equation}
where $V_0 = V_{\mathrm{prod}}$ is the raw product map and $V_2$ is built from the pair-creation operator $h^\dagger$~\cite{GarjaniArdonne2016} applied on the fine lattice.

\subsubsection{The deficit}

The Gram matrix $G = V_0^\dagger V_0$ is diagonal with entries $g(\vec{x})^2 \le 1$. The deficit
\begin{equation}
    D = I - G = \operatorname{diag}\bigl(1 - g(\vec{x})^2\bigr)
\end{equation}
vanishes for $N = 0, 1$ (no hard-core collisions possible) and grows with the number of adjacent particle pairs. For D4 wavelets, the deficit values are approximately
\begin{center}
\begin{tabular}{cc}
$N$ & max deficit \\
\hline
0, 1 & 0 \\
2 & 0.029 \\
3 & 0.057 \\
4 & 0.083
\end{tabular}
\end{center}
These values are \emph{category-independent}: Fibonacci and Ising produce identical deficit spectra.

\subsubsection{Composition with pair creation}

The pair-creation operator $h^\dagger$ on the fine lattice maps $\catH_N(2L) \to \catH_{N+2}(2L)$ by creating vacuum pairs at adjacent vacant sites~\cite{GarjaniArdonne2016}. The raw composition $h^\dagger V_0$ maps $\catH_N(L) \to \catH_{N+2}(2L)$, spreading coarse particles via $R$ and then inserting pairs at the newly created vacant fine sites.

The naive construction $V = V_0 + \alpha\, h^\dagger V_0$ fails for two reasons:
\begin{enumerate}
    \item \textbf{Cross-sector overlap}: $V_0$ also maps coarse $(N{+}2)$-particle states into the fine $(N{+}2)$-particle sector, so $V_0^\dagger (h^\dagger V_0) \neq 0$.
    \item \textbf{Column interference}: the $h^\dagger V_0$ columns for different coarse states are not orthogonal (different coarse states can produce the same fine state via pair creation at different bonds).
\end{enumerate}

\subsubsection{Three-step orthogonalisation}

Both obstructions are resolved by a systematic orthogonalisation procedure:

\medskip\noindent
\textbf{Step 1: Projection.} Let $\Pi_{V_0} = V_0 G^{-1} V_0^\dagger$ be the projector onto $\operatorname{Im}(V_0)$, where $G^{-1} = \operatorname{diag}(1/g(\vec{x})^2)$ (well-defined since $g > 0$ for all occupied states). Define
\begin{equation}
    V_2^{\perp} = (I - \Pi_{V_0})\, h^\dagger V_0.
\end{equation}
This ensures $V_0^\dagger V_2^{\perp} = 0$ exactly.

\medskip\noindent
\textbf{Step 2: L\"owdin orthogonalisation.} Restrict to the columns $J = \{j : D_{jj} > 0\}$ with nonzero deficit. The Gram matrix of the projected columns,
\begin{equation}
    G_2 = (V_2^{\perp})_J^\dagger (V_2^{\perp})_J,
\end{equation}
is positive definite (verified numerically). Apply symmetric orthogonalisation:
\begin{equation}
    (V_2^{\perp})_J \;\longmapsto\; (V_2^{\perp})_J \cdot G_2^{-1/2}.
\end{equation}
This gives orthonormal $V_2$ columns within the complement of $\operatorname{Im}(V_0)$.

\medskip\noindent
\textbf{Step 3: Scaling.} Multiply each column $j \in J$ by $\sqrt{D_{jj}}$ to fill the deficit:
\begin{equation}
\label{eq:V2-formula}
    V_2 = (I - \Pi_{V_0})\, h^\dagger V_0 \cdot G_2^{-1/2} \cdot D_J^{1/2}.
\end{equation}

\begin{proposition}[Number-changing isometry]
\label{prop:number-changing}
The map $V = V_0 + V_2$ with $V_2$ given by~\eqref{eq:V2-formula} satisfies $V^\dagger V = I$.
\end{proposition}

\begin{proof}
By construction:
\begin{equation}
    V^\dagger V = V_0^\dagger V_0 + \underbrace{V_0^\dagger V_2}_{= 0} + \underbrace{V_2^\dagger V_0}_{= 0} + V_2^\dagger V_2 = G + D = I.
\end{equation}
The cross-terms vanish by Step~1. The diagonal term $V_2^\dagger V_2 = D$ by Steps~2--3: the L\"owdin orthogonalisation makes the columns orthonormal, and the scaling adjusts each to have squared norm $D_{jj}$.
\end{proof}

\begin{remark}
The construction requires $G_2$ to be positive definite on the deficit subspace, i.e., the pair-creation operator must generate enough independent fine-state components orthogonal to $\operatorname{Im}(V_0)$. This has been verified numerically for Fibonacci and Ising at $L = 3, 4$ with D4 wavelets.
\end{remark}

\subsubsection{Numerical results}

We have verified $\|V^\dagger V - I\| < 10^{-14}$ (machine precision) for:
\begin{itemize}
    \item Fibonacci category, $L = 3, 4$, D4 wavelet
    \item Ising category (multi-species: $\sigma, \psi$), $L = 3, 4$, D4 wavelet
\end{itemize}

\noindent
For the Hamiltonian intertwining $V^\dagger H_{2L} V \approx a H_L + b I$, the number-changing isometry gives slightly better residuals than the normalized product map:

\begin{center}
\begin{tabular}{lccc}
& $L$ & Normalized product & Number-changing \\
\hline
Fibonacci & 3 & 0.448 & 0.434 \\
Fibonacci & 4 & 0.393 & 0.378
\end{tabular}
\end{center}

\subsubsection{Physical interpretation}

The number-changing isometry realises the fine-graining as a two-step physical process:
\begin{enumerate}
    \item \textbf{Spread}: each coarse particle spreads to fine sites via the wavelet filter $R$ (preserving $N$).
    \item \textbf{Pair creation}: vacuum pairs are created at vacant fine sites to fill the deficit left by hard-core collisions (changing $N \to N+2$).
\end{enumerate}
The pair-creation contribution is small (scaling factors $\sim \sqrt{D/\|h^\dagger V_0\|^2} \sim 0.1$) but structurally necessary. This provides the first concrete evidence that Garjani--Ardonne pair creation~\cite{GarjaniArdonne2016} plays a role in the OAR continuum limit~\cite{Stottmeister2022}, not just in the lattice Hamiltonian.

\subsection{Hierarchy of constructions}

Table~\ref{tab:hierarchy} summarises which fine-graining constructions are valid in each regime.

\begin{table}[h]
\centering
\begin{tabular}{lccc}
\hline
\textbf{Construction} & \textbf{Single species} & \textbf{Multi-species} & \textbf{Number-preserving} \\
\hline
Product map (Haar) & \checkmark & \checkmark & yes \\
Product map (D4) & $\times$ & $\times$ & yes \\
Slater $\Lambda^N(R) \otimes I$ & \checkmark & $\times$ & yes \\
Normalized product & \checkmark & \checkmark & yes \\
Number-changing~\eqref{eq:number-changing} & \checkmark & \checkmark & no \\
\hline
\end{tabular}
\caption{Fine-graining isometry constructions. The normalized product map~\eqref{eq:normalized-product} and number-changing isometry~\eqref{eq:number-changing} are both universal. The former is simpler; the latter is physically motivated and gives slightly better Hamiltonian intertwining.}
\label{tab:hierarchy}
\end{table}

\subsection{Multi-step composition and the inductive limit}

The fine-graining isometries compose correctly: $V_{L \to 4L} = V_{2L \to 4L} \circ V_{L \to 2L}$ is isometric whenever each factor is. This is verified numerically for sVec, Fibonacci, and Ising with D4 wavelets at $L = 3$.

The sequence of lattice algebras $\mathcal{A}_L \xrightarrow{\alpha_{2L,L}} \mathcal{A}_{2L} \xrightarrow{\alpha_{4L,2L}} \mathcal{A}_{4L} \to \cdots$ under the fine-graining $*$-homomorphisms $\alpha_{2L,L}(A) = V^\dagger A V$ forms a directed system whose inductive limit $\mathcal{A}_\infty = \varinjlim \mathcal{A}_N$ is the continuum C*-algebra of observables. The existence and uniqueness of this limit is guaranteed by the general theory of inductive limits of C*-algebras~\cite{OsborneStottmeister2022}.

\subsection{Categorical interpretation}

The key insight is that the \emph{position} and \emph{internal} degrees of freedom play fundamentally different roles:
\begin{itemize}
    \item \textbf{Position space}: For hard-core particles, the $N$-body position space is $\binom{L}{N}$-dimensional, isomorphic to $\Lambda^N(\mathbb{C}^L)$ as a vector space (though the isomorphism involves a choice of ordering). The fine-graining map $R$ acts only on positions.
    \item \textbf{Internal space}: The fusion-tree space $\Mor(X_c, X_{a_1} \otimes \cdots \otimes X_{a_N})$ encodes the categorical data. It is independent of positions and is preserved by the fine-graining map.
\end{itemize}

The categorical data (F-symbols, R-symbols, quantum dimensions) enters the physics through the \emph{Hamiltonian} --- hopping, interaction, braiding, and pair creation terms --- not through the fine-graining map itself. The fine-graining is a purely geometric operation on the lattice, while the anyonic structure lives in the dynamics.
